\documentclass{article}
\usepackage[utf8]{inputenc}
\usepackage{a4wide}

\begin{document}

\section*{สายสืบ}

วริยาได้รับงานหนึ่งจากผู้จ้างวานให้ปลอมตัวเข้าไปทำงานในบริษัทยักษ์ใหญ่เจ้าหนึ่งที่ให้บริการแอพลิเคชันด้าน

social networks สิ่งที่ผู้ว่าจ้างต้องการคือให้ดูว่ามี user ใดบ้างที่มีโอกาสเป็นคนคนเดียวกัน



แอพลิเคชัน social networks นี้มี n account เรียงลำดับตั้งแต่ 1 ถึง n ซึ่งบาง account

ก็เป็นเพื่อนกัน (หาก I เป็นเพื่อนกับ j ถือว่า j เป็นเพื่อนกับ I เช่นเดียวกัน)



คู่ account (i, j) จะถูกเรียกว่าเป็น "คู่เหมือน" ในกรณีที่ สำหรับ account k ใดๆ ที่ไม่ใช่

i หรือ j จะต้องเป็นไปตามเงื่อนไขอย่างใดอย่างหนึ่ง:



\begin{enumerate}

\tightlist

\item

  โปรไฟล์ k เป็นเพื่อนกับทั้ง i และ j

\item

  โปรไฟล์ k ไม่ได้เป็นเพื่อนกับทั้ง i และ j

\end{enumerate}



โดยที่ i และ j จะเป็นเพื่อนกันหรือไม่ก็ได้



\textbf{งานของคุณ}-- ช่วยวริยาคำนวณจำนวนคู่ account ที่เป็น "คู่เหมือน" โดยที่ (a, b)

กับ (b, a) ถือว่าเป็นคู่เดียวกัน



{\def\LTcaptype{none} % do not increment counter

\begin{longtable}[]{@{}

  >{\raggedright\arraybackslash}p{(\linewidth - 4\tabcolsep) * \real{0.3333}}

  >{\raggedright\arraybackslash}p{(\linewidth - 4\tabcolsep) * \real{0.3333}}

  >{\raggedright\arraybackslash}p{(\linewidth - 4\tabcolsep) * \real{0.3333}}@{}}

\toprule\noalign{}

\begin{minipage}[b]{\linewidth}\centering

input

\end{minipage} & \begin{minipage}[b]{\linewidth}\centering

output

\end{minipage} & \begin{minipage}[b]{\linewidth}\centering

comment

\end{minipage} \\

\midrule\noalign{}

\endhead

\bottomrule\noalign{}

\endlastfoot

\begin{minipage}[t]{\linewidth}\raggedright

3 3\\

1 2\\

2 3\\

1 3\\

\strut

\end{minipage} & 3 & \begin{minipage}[t]{\linewidth}\raggedright

ในกรณีนี้ ทุก account มีความสัมพันธ์กับ account อื่น ๆ ซึ่งสามารถจับคู่ account ได้ทั้งหมด

3 คู่ ได้แก่ (1, 2), (1, 3), (2, 3)\\

\strut

\end{minipage} \\

3 0 & 3 & \begin{minipage}[t]{\linewidth}\raggedright

ในกรณีนี้ ไม่มีความสัมพันธ์เพื่อนใด ๆ เลย ดังนั้นทุกคู่ account จะเป็น "คู่เหมือน" (เพราะไม่มี

account ไหนที่เป็นเพื่อนกัน)\\

\strut

\end{minipage} \\

\begin{minipage}[t]{\linewidth}\raggedright

4 1\\

1 3\\

\strut

\end{minipage} & 2 & \begin{minipage}[t]{\linewidth}\raggedright

ในกรณีนี้มีเพียงแค่ความสัมพันธ์เพื่อนระหว่าง account 1 และ 3 ดังนั้นคู่ account ที่เป็น

"คู่เหมือน" ได้แก่ (1, 3) และ (2, 4)\\

\strut

\end{minipage} \\

\end{longtable}

}



\subsection*{Input}
บรรทัดที่ 1: จำนวนเต็ม n (จำนวน account ทั้งหมด 1 ≤ n ≤ 106) และ m

(จำนวนความพันธ์ที่เป็นเพื่อนกัน 1 ≤ m ≤ 106) แยกด้วยช่องว่าง



บรรทัดที่ 2 เป็นต้นไป จำนวน m บรรทัด: ความสัมพันธ์ที่เป็นเพื่อนกันในรูปแบบ v u โดยที่

1 ≤ v, u ≤ n, v ≠ u ซึ่งหมายถึง account v และ account u เป็นเพื่อนกัน

ต้องมั่นใจว่าใน input มีแค่คู่ v u หรือ u v แค่อย่างใดอย่างหนึ่ง และ ไม่มี v v หรือ u u



\subsection*{Output}
บรรทัดที่ 1: จำนวนคู่ account ที่เป็น "คู่เหมือน"\\



\end{document}
